% Part: propositional-logic
% Chapter: syntax-and-semantics
% Section: semantic-notions

\documentclass[../../../include/open-logic-section]{subfiles}

\begin{document}

\olfileid[fr]{pl}{syn}{sem}

\olsection{Notions sémantiques}

Définissons les notions sémantiques suivantes :

\begin{defn}
\begin{enumerate}
\item Une !!{formula}~$!A$ est \emph{satisfaisable} s'il existe
  une~$\pAssign{v}$ telle que $\pSat{v}{!A}$ ; elle est
  \emph{insatisfaisable} s'il n'existe aucune $\pAssign{v}$ telle
  que $\pSat{v}{!A}$.
\item Une !!{formula}~$!A$ est une \emph{tautologie} si
  $\pSat{v}{!A}$ pour toute !!{valuation}~$\pAssign{v}$.
\item Une !!{formula}~$!A$ est \emph{contingente} si elle est
  satisfaisable sans être une tautologie.
\item Si $\Gamma$ est un ensemble de !!{formula}s, on écrit
  $\Gamma \Entails !A$ (« $!A$ est conséquence logique de $\Gamma$ »)
  si et seulement si $\pSat{v}{!A}$ pour toute
  !!{valuation}~$\pAssign{v}$ telle que $\pSat{v}{\Gamma}$.
\item Si $\Gamma$ est un ensemble de !!{formula}s, $\Gamma$ est
  \emph{satisfaisable} s'il existe une !!{valuation}~$\pAssign{v}$
  telle que $\pSat{v}{\Gamma}$ ; il est \emph{insatisfaisable}
  dans le cas contraire.
\end{enumerate}
\end{defn}

\begin{prob}
Pour chacune des quatre !!{formula}s suivantes, déterminer si elle est
(a)~satisfaisable, (b)~une tautologie et (c)~contingente.
\begin{enumerate}
  \item \((\Obj p_0 \lif (\lnot \Obj p_1 \lif \lnot \Obj p_0))\).
  \item \(((\Obj p_0 \land \lnot \Obj p_1) \lif (\lnot \Obj p_0 \land \Obj p_2)) \liff ((\Obj p_2 \lif \Obj p_0) \lif (\Obj p_0 \lif \Obj p_1))\).
  \item \((\Obj p_0 \liff \Obj p_1) \lif (\Obj p_2 \liff \lnot \Obj p_1)\).
  \item \(((\Obj p_0 \liff (\lnot \Obj p_1 \land \Obj p_2)) \lor (\Obj p_2 \lif (\Obj p_0 \liff \Obj p_1)))\).
\end{enumerate}
\end{prob}

\begin{prop}
\ollabel{prop:semanticalfacts}
\begin{enumerate}
\item $!A$ est une tautologie si et seulement si
  $\emptyset \Entails !A$.
\item Si $\Gamma \Entails !A$ et $\Gamma \Entails !A \lif !B$,
  alors $\Gamma \Entails !B$.
\item Si $\Gamma$ est satisfaisable, chacun de ses sous-ensembles
  finis l'est aussi.
\item \ollabel{def:monotonicity} Monotonie : si
  $\Gamma \subseteq \Delta$ et $\Gamma \Entails !A$, alors
  $\Delta \Entails !A$.
\item \ollabel{def:Cut} Transitivité : si $\Gamma \Entails !A$ et
  $\Delta \cup \{!A\} \Entails !B$, alors
  $\Gamma \cup \Delta \Entails !B$.
\end{enumerate}
\end{prop}

\begin{proof}
La démonstration est laissée en exercice.
\end{proof}

\begin{prob}
Démontrer la \olref[pl][syn][sem]{prop:semanticalfacts}.
\end{prob}

\begin{prop}\ollabel{prop:entails-unsat}
  $\Gamma \Entails !A$ si et seulement si
  $\Gamma \cup \{\lnot !A\}$ est insatisfaisable.
\end{prop}

\begin{proof}
La démonstration est laissée en exercice.
\end{proof}

\begin{prob}
Démontrer la \olref[pl][syn][sem]{prop:entails-unsat}.
\end{prob}

\begin{thm}[Théorème sémantique de déduction]
  \ollabel{thm:sem-deduction}
  $\Gamma \Entails !A \lif !B$ si et seulement si
  $\Gamma \cup \{!A\} \Entails !B$.
\end{thm}

\begin{proof}
La démonstration est laissée en exercice.
\end{proof}

\begin{prob}
Démontrer le \olref[pl][syn][sem]{thm:sem-deduction}.
\end{prob}

\end{document}
